\documentclass[12pt,a4paper]{report}

% ============================================================
%  IEEE-Compliant FYP Report
%  Formatting per Project_Guidelines.pdf & latex-ieee-formatting skill
% ============================================================

% --- Packages ---
\usepackage{fontspec}
\setmainfont{Simplified Arabic}

\usepackage{tikz}
\usetikzlibrary{positioning, shapes.geometric, arrows.meta, backgrounds, fit, calc}
\usepackage[left=1.5in,top=1in,bottom=1in,right=1in]{geometry} % 1.5in binding margin
\usepackage{setspace}
\usepackage{titlesec}
\usepackage[table]{xcolor}
\usepackage{booktabs}
\usepackage{longtable}
\usepackage{array}
\usepackage{tabularx}
\usepackage{hyperref}
\usepackage{amsmath, amssymb}       % For ML formulas
\usepackage{graphicx}
\usepackage[nottoc,notlof,notlot]{tocbibind}
\usepackage{parskip}
\usepackage{caption}
\usepackage{url}
\usepackage{fancyhdr}
\usepackage{listings}
\usepackage{algorithm2e}
\usepackage{polyglossia}
\setdefaultlanguage{english}
\setotherlanguage{arabic}
\newfontfamily\arabicfont[Script=Arabic, Ligatures=Required]{Simplified Arabic}

\usepackage[backend=biber]{biblatex} % For citations
\usepackage{minted}                 % For your Go/Laravel code snippets
\usepackage{neuralnetwork}
\usepackage{verbatim}
\usepackage{alltt}

% --- Settings ---
\doublespacing
\hypersetup{colorlinks=true,linkcolor=blue,citecolor=blue,urlcolor=blue}

% Chapter title formatting: CHAPTER X prefix, centered, bold
\titleformat{\chapter}[display]
  {\normalfont\huge\bfseries\centering}
  {CHAPTER \thechapter}{0pt}{\Huge}
\titleformat{\section}{\normalfont\Large\bfseries}{\thesection}{1em}{}
\titleformat{\subsection}{\normalfont\large\bfseries}{\thesubsection}{1em}{}

% ============================================================
\begin{document}

% ==============================
%  PART 1: FRONT MATTER
%  Roman numeral page numbering
% ==============================
\pagenumbering{roman}
\titlespacing*{\chapter}{0pt}{-30pt}{20pt}

% -------------------- TITLE PAGE --------------------
\begin{titlepage}

    % Header Table
    \noindent\makebox[\textwidth]{
        \small
        \begin{tabular}{c c c}
        \begin{tabular}[c]{@{}l@{}}
            \textbf{REPUBLIC OF YEMEN} \\
            \textbf{Ministry of Higher Education of Sciences} \\
            \textbf{TAIZ UNIVERSITY} \\
            \textbf{Faculty Of Engineering IT} \\
            \textbf{SE Department}
        \end{tabular} 
        & 
        \begin{tabular}[c]{@{}c@{}}
            \vspace{0.1cm}
            \resizebox{2.5cm}{!}{
                \definecolor{logoblue}{RGB}{0, 130, 235}
                \definecolor{logolightblue}{RGB}{120, 190, 255}
                \definecolor{logoyellow}{RGB}{255, 235, 20}
                \definecolor{atomblue}{RGB}{140, 200, 220}
                \begin{tikzpicture}[line cap=round, line join=round]
                    % ==========================================
                    % 1. BOTTOM BANNER
                    % ==========================================
                    \draw [line width=3.5pt, fill=logoyellow] (-5, 0.7) rectangle (5, -1.8);
                    
                    % Banner inner horizontal line
                    \draw [line width=1.5pt] (-4.8, -0.45) -- (4.8, -0.45);
                    
                    % Banner Texts
                    \node[fill=logoyellow, inner xsep=10pt, inner ysep=2pt] at (0, -0.45) {\huge \textbf{\textarabic{جامعة تعز}}};
                    \node at (0, -1.25) {\Large \rmfamily TAIZ UNIVERSITY};
                    
                    % Dates
                    \node[fill=logoyellow, inner sep=3pt] at (-4.2, -0.45) {\normalsize 1993\textarabic{م}};
                    \node[fill=logoyellow, inner sep=3pt] at (4.2, -0.45) {\normalsize 1414\textarabic{هـ}};

                    % ==========================================
                    % 2. MAIN ARCH BACKGROUND
                    % ==========================================
                    \draw[line width=1.5pt, draw=black, top color=logolightblue, bottom color=logoblue] 
                        (-5, 1.1) -- (-5, 4) 
                        .. controls (-5, 8.5) and (-2.5, 11.5) .. (0, 12.5) 
                        .. controls (2.5, 11.5) and (5, 8.5) .. (5, 4) 
                        -- (5, 1.1) -- cycle;

                    % ==========================================
                    % 3. SUN AND RAYS
                    % ==========================================
                    \draw[line width=2pt, fill=logoyellow] (0, 9.5) circle (2.0);
                    
                    \foreach \a/\len in {
                        218/1.5, 224.5/0.75, 231/1.5, 237.5/0.75, 244/1.5, 250.5/0.75, 257/1.5, 263.5/0.75, 
                        270/1.5, 
                        276.5/0.75, 283/1.5, 289.5/0.75, 296/1.5, 302.5/0.75, 309/1.5, 315.5/0.75, 322/1.5
                    } {
                        \draw [logoyellow, line width=2pt] (0, 9.5) ++(\a:2.3) -- ++(\a:\len);
                    }
                    
                    \node at (0, 8.8) {\large \textbf{\textarabic{وقل ربي زدني علما}}};

                    % ==========================================
                    % 4. OPEN BOOK (Inside the Sun)
                    % ==========================================
                    \begin{scope}[shift={(0, 9.8)}]
                        \fill [black] (-1.3, -0.4) .. controls (-0.7, -0.3) .. (0, -0.6) 
                                   .. controls (0.7, -0.3) .. (1.3, -0.4) -- (1.2, 0) -- (0, -0.2) -- (-1.2, 0) -- cycle;
                        
                        \draw[line width=1.5pt, fill=white] (0, -0.5) 
                            .. controls (-0.5, -0.2) and (-1.0, -0.3) .. (-1.3, -0.3) 
                            -- (-0.9, 0.6) 
                            .. controls (-0.6, 0.6) and (-0.2, 0.7) .. (0, 0.4) -- cycle;
                            
                        \draw[line width=1.5pt, fill=white] (0, -0.5) 
                            .. controls (0.5, -0.2) and (1.0, -0.3) .. (1.3, -0.3) 
                            -- (0.9, 0.6) 
                            .. controls (0.6, 0.6) and (0.2, 0.7) .. (0, 0.4) -- cycle;
                            
                        \draw [line width=1.2pt] (-0.2, 0.35) -- (-0.4, 0.4); 
                        \draw[line width=1.2pt] (-0.5, 0.42) -- (-0.8, 0.5);
                        \draw[line width=1.2pt] (-0.2, 0.1) -- (-0.5, 0.15); 
                        \draw[line width=1.2pt] (-0.6, 0.17) -- (-0.9, 0.25);
                        \draw[line width=1.2pt] (-0.2, -0.15) -- (-0.6, -0.1); 
                        \draw[line width=1.2pt] (-0.7, -0.08) -- (-1.0, 0.0);
                        
                        \draw[line width=1.2pt] (0.2, 0.35) -- (0.4, 0.4); 
                        \draw[line width=1.2pt] (0.5, 0.42) -- (0.8, 0.5);
                        \draw[line width=1.2pt] (0.2, 0.1) -- (0.5, 0.15); 
                        \draw[line width=1.2pt] (0.6, 0.17) -- (0.9, 0.25);
                        \draw[line width=1.2pt] (0.2, -0.15) -- (0.6, -0.1); 
                        \draw[line width=1.2pt] (0.7, -0.08) -- (1.0, 0.0);
                    \end{scope}

                    % ==========================================
                    % 5. ATOM SYMBOL (Center)
                    % ==========================================
                    \begin{scope}[shift={(0, 4.0)}, scale=0.7, line width=1pt]
                        \foreach \angle in {0, 45, 90, 135} {
                            \fill [atomblue, rotate=\angle] (0,0) ellipse (0.4 and 1.5);
                        }
                        \foreach \angle in {0, 45, 90, 135} {
                            \draw [rotate=\angle] (0,0) ellipse (0.4 and 1.5);
                        }
                        \draw [fill=logoyellow] (0,0) circle (0.2);
                        \fill [black] (0,0) circle (0.06);
                    \end{scope}

                    % ==========================================
                    % 6. BUILDINGS & MINARET (Bottom Layer)
                    % ==========================================
                    \draw[line width=6pt] (-4.6, 1.5) -- (4.6, 1.5);

                    \begin{scope}[shift={(-3.2, 1.5)}, line width=1.5pt]
                        \draw [fill=white] (-0.7, 0) rectangle (0.7, 0.3);
                        \draw[fill=white] (-0.4, 0.3) rectangle (0.4, 2.5);
                        \draw [fill=white] (-0.4, 2.5) -- (-0.6, 3.0) -- (0.6, 3.0) -- (0.4, 2.5) -- cycle;
                        \draw [fill=white] (-0.6, 3.0) rectangle (0.6, 3.3);
                        \draw [fill=white] (-0.25, 3.3) rectangle (0.25, 4.2);
                        \draw [fill=white] (-0.25, 4.2) -- (-0.35, 4.5) -- (0.35, 4.5) -- (0.25, 4.2) -- cycle;
                        \draw [fill=white] (-0.25, 4.5) rectangle (0.25, 4.8);
                        \draw[fill=white] (-0.25, 4.8) -- (0, 5.5) -- (0.25, 4.8) -- cycle;
                    \end{scope}

                    \draw[line width=1.5pt, fill=white] 
                        (-2.4, 1.5) -- (-2.4, 2.0)
                        -- (-1.875, 2.7) -- (-1.35, 2.0)
                        -- (-0.825, 2.7) -- (-0.3, 2.0)
                        -- (0.225, 2.7)  -- (0.75, 2.0)
                        -- (1.275, 2.7)  -- (1.8, 2.0)
                        -- (2.325, 2.7)  -- (2.85, 2.0)
                        -- (3.375, 2.7)  -- (3.9, 2.0)
                        -- (3.9, 1.5) -- cycle;
                        
                    \foreach \x in {-1.875, -0.825, 0.225, 1.275, 2.325, 3.375} {
                        \fill [black] (\x, 1.8) circle (0.12);
                    }
                \end{tikzpicture}
            }
            \vspace{0.1cm}
        \end{tabular} 
        & 
        \begin{tabular}[c]{@{}r@{}}
            \textbf{\textarabic{الجمهورية اليمنية}} \\
            \textbf{\textarabic{وزارة التعليم العالي والبحث العلمي}} \\
            \textbf{\textarabic{جامعة تعز}} \\
            \textbf{\textarabic{كلية الهندسة وتقنية المعلومات}} \\
            \textbf{\textarabic{قسم البرمجيات}}
        \end{tabular} \\
        \end{tabular}
    }

    \vspace{0.8cm}

    \begin{center}
        % Title
        \textbf{\Large Design and Development of a Unified Gesture Framework} \\[0.2cm]
        \textbf{\Large for Cross-Platform Mobile Applications}
        
        \vspace{0.6cm}
        
        % Subtitle
        \textbf{\normalsize A Final Year Project Report} \\[0.5cm]
        \normalsize Submitted in partial fulfillment of the requirements for the Degree of \\
        Bachelor of Science in Software Engineering
        
        \vspace{0.6cm}
        
        % Team
        \textbf{\normalsize Prepared by:} \\[0.3cm]
        \normalsize - Ziad Faisal Gazem \\[0.1cm]
        \normalsize - Gamal Sadek Saeed \\[0.1cm]
        \normalsize - Hafeed Shaheed \\[0.1cm]
        \normalsize - Moharram Al-jaradi \\[0.1cm]
        
        \vspace{0.6cm}
        
        % Supervisor
        \textbf{\normalsize Supervisor} \\[0.3cm]
        \normalsize Dr. Ziad Al-Robih
        
        \vfill
        
        % Footer
        \textbf{\normalsize College of Engineering and Information Technology} \\[0.2cm]
        \textbf{\normalsize Taiz University} \\[0.2cm]
        \textbf{\normalsize Taiz, Yemen} \\[0.2cm]
        \textbf{\normalsize }
    \end{center}

\end{titlepage}


% -------------------- DECLARATION --------------------
\begin{Arabic}
\chapter*{الإقرار}

نحن، الموقعون أدناه، طلاب كلية الهندسة وتقنية المعلومات في جامعة تعز، نقر بأن مشروع التخرج المعنون ``مكتبة إيماءات اللمس المتقدمة لتطبيقات React Native'' هو عملنا الأصلي.

وحسب علمنا، فإن هذا المستند لا يحتوي على أي مواد سبق نشرها أو كتابتها من قبل شخص آخر، ولا مواد تم قبولها للحصول على أي درجة أو دبلوم آخر في هذه الجامعة أو أي مؤسسة تعليمية أخرى، باستثناء ما تمت الإشارة إليه في النص. ونقر باستخدام أطر العمل مفتوحة المصدر والمراجع التقنية وفقاً لتراخيصها، والتي تم توثيقها بالكامل داخل التقرير.

% -------------------- APPROVAL PAGE --------------------
\chapter*{صفحة الاعتماد}

تم فحص واعتماد مشروع التخرج هذا المعنون ``مكتبة إيماءات اللمس المتقدمة لتطبيقات React Native'' كجزء من متطلبات الحصول على درجة البكالوريوس في هندسة البرمجيات.

\vspace{1.5cm}

\noindent\textbf{مشرف المشروع:}\\
\noindent\rule{8cm}{0.4pt}\\
د. زياد الربيع\\
التاريخ: \rule{4cm}{0.4pt}

\vspace{0.75cm}

\noindent\textbf{الممتحن الداخلي:}\\
\noindent\rule{8cm}{0.4pt}\\
الاسم: \rule{5cm}{0.4pt}\\
التاريخ: \rule{4cm}{0.4pt}

\vspace{0.75cm}

\noindent\textbf{رئيس القسم:}\\
\noindent\rule{8cm}{0.4pt}\\
الاسم: \rule{5cm}{0.4pt}\\
التاريخ: \rule{4cm}{0.4pt}

% -------------------- DEDICATION --------------------
\chapter*{الإهداء}

نهدي هذا المشروع إلى عائلاتنا، التي لولا دعمها اللامحدود، وصبرها، وتشجيعها لما أمكننا إتمام مسيرتنا الأكاديمية. كما نهدي هذا العمل إلى أساتذتنا وإلى مجتمع البرمجيات مفتوحة المصدر، الذين ألهمتنا أدواتهم الأساسية ومعارفهم المشتركة. وأخيراً، نهدي هذا العمل إلى وطننا الحبيب اليمن، آملين أن نساهم في تقدمه التكنولوجي ومستقبله المشرق.

% -------------------- ACKNOWLEDGEMENTS --------------------
\chapter*{شكر وتقدير}

أولاً وقبل كل شيء، نشكر الله عز وجل على منحنا القوة والصحة والحكمة لإتمام هذا المشروع بنجاح.

نود أن نعرب عن أعمق امتناننا لمشرف مشروعنا، د. زياد الربيع، لتوجيهاته القيمة، وملاحظاته البناءة، ودعمه المستمر طوال فترة تطوير هذا البحث.

كما نتقدم بخالص الشكر والتقدير لأعضاء هيئة التدريس في كلية الهندسة وتقنية المعلومات بجامعة تعز لتزويدنا بالمعرفة الأساسية اللازمة للقيام بمشروع بهذا الحجم التقني.
\end{Arabic}


% -------------------- ARABIC ABSTRACT --------------------
\chapter*{\textarabic{الملخص}}

\begin{Arabic}
\begingroup
\fontsize{11pt}{13.2pt}\selectfont
مع تطور تطبيقات الهواتف الذكية وتزايد تعقيد واجهات المستخدم، أصبحت إيماءات اللمس عنصراً حيوياً لتوفير تجربة مستخدم سلسة وتفاعلية. يواجه مطورو التطبيقات متعددة المنصات (مثل \textenglish{React Native}) تحديات مستمرة تتمثل في تأخر الاستجابة الملحوظ الناتج عن تمرير الأحداث عبر الجسر (\textenglish{The Bridge}) بين الواجهة الأصلية للنظام وبيئة \textenglish{JavaScript}. يهدف هذا المشروع إلى تصميم وتطوير إطار عمل موحد ومتقدم (\textenglish{Gesture Framework}) مخصص لمعالجة إيماءات اللمس بمختلف مستوياتها، بحيث يتجاوز هذه القيود المعمارية ويضمن أداءً فائقاً بمعدل 60 إطاراً في الثانية.

يُقدم هذا الإطار منظومة شاملة تتكون من خمس حزم أساسية تغطي طيفاً واسعاً من التفاعلات: بدءاً من الإيماءات الأساسية والمنفصلة (كالضغطة المزدوجة والضغط المطول)، مروراً بإيماءات السحب والانزلاق المتصلة (\textenglish{Drag/Pan})، وصولاً إلى الإيماءات متعددة الأصابع والتحويلات الهندسية المعقدة (مثل إيماءة القرص للتكبير والتدوير المتزامن). ويعتمد التصميم المعماري للمشروع بشكل جوهري على نقل معالجة الإيماءات بالكامل إلى مسار الواجهة الأصلية (\textenglish{UI Thread}) بالاعتماد على مكتبة \textenglish{react-native-gesture-handler} والمتغيرات المشتركة في \textenglish{Reanimated}، مما يُلغي الاختناقات المرتبطة بالجسر. بالإضافة إلى ذلك، يوسع الإطار قدراته ليدمج بيانات الحساسات (مثل مستشعر التقارب \textenglish{Proximity}) ومحرك رسم عالي الأداء لتنفيذ تطبيقات التوقيع الرقمي والتفاعلات المتقدمة.

أظهرت التقييمات المعمارية نجاح هذا النهج في تحقيق تفاعل شبه فوري وتوفير استقرار مرجعي عالٍ. ومن خلال تقديم واجهة برمجة تطبيقات (\textenglish{API}) نظيفة ومستقرة عبر خطافات (\textenglish{Hooks}) مخصصة، يوفر الإطار حلاً جذرياً يُسهّل على مجتمع المطورين دمج التفاعلات الحركية المعقدة في تطبيقاتهم، مما يعزز من جودة تجربة المستخدم عبر مختلف المنصات دون المساس بسرعة الاستجابة.
\par\endgroup
\end{Arabic}


% -------------------- ENGLISH ABSTRACT --------------------
\chapter*{ABSTRACT}

\begingroup
\fontsize{11pt}{13.2pt}\selectfont
As smartphones evolve and user interfaces grow increasingly complex, touch gestures have become a vital component for delivering seamless and interactive user experiences. Cross-platform application developers (such as those using React Native) face persistent challenges regarding noticeable response latency caused by event serialization across ``The Bridge'' between the native platform and the JavaScript environment. This project aims to design and develop a unified, advanced Gesture Framework tailored for processing touch gestures across all complexities, bypassing architectural constraints to guarantee high performance at a consistent 60 frames per second (fps).

The framework introduces a comprehensive ecosystem comprising five primary packages that cover a wide spectrum of interactions. These range from discrete, basic touch events (e.g., double tap, long press), to continuous swipe and drag interactions (Drag/Pan), up to multi-finger combinations and complex geometric transformations (such as simultaneous pinch-to-zoom and rotation). The core architectural design of the project relies on migrating gesture processing entirely to the Native UI Thread. By leveraging \textit{react-native-gesture-handler} alongside shared values from \textit{Reanimated}, the framework eliminates bridge-related bottlenecks. Furthermore, the architecture extends its capabilities to integrate hardware sensor data (e.g., proximity sensors) and a high-performance drawing engine tailored for digital signature applications and advanced continuous interactions.

Architectural evaluations demonstrate the success of this approach in achieving near-instantaneous interaction and providing high referential stability. By delivering a clean and stable Application Programming Interface (API) through custom React Hooks, the framework offers a radical solution that facilitates the integration of complex kinematic interactions into the broader developer community's applications. This ultimately enhances the quality of user experiences across different platforms without compromising responsiveness.
\par\endgroup

% -------------------- TOC / LOF / LOT --------------------
\listoftables
\listoffigures
\tableofcontents

% -------------------- LIST OF ABBREVIATIONS --------------------
\begin{Arabic}
\chapter*{قائمة الاختصارات}
\end{Arabic}

\begin{tabular}{@{}ll@{}}
\textbf{UI/UX} & User Interface / User Experience \\
\textbf{HAL}   & Hardware Abstraction Layer \\
\textbf{API}   & Application Programming Interface \\
\textbf{NPM}   & Node Package Manager \\
\end{tabular}

% -------------------- GLOSSARY --------------------
\begin{Arabic}
\chapter*{مسرد المصطلحات}

\begin{description}
\item[واجهة برمجة التطبيقات (API)] مجموعة من القواعد والبروتوكولات التي تسمح لتطبيقات البرامج المختلفة بالتواصل مع بعضها البعض.

\item[الجسر (The Bridge)] آلية في React Native تسمح لبيئة JavaScript بالتواصل مع الواجهات الأصلية للنظام لمعالجة الإيماءات.

\item[الشاشة السعوية (Capacitive Screen)] تقنية شاشات اللمس الحديثة التي تعتمد على الخصائص الكهربائية لجسم الإنسان لاستشعار اللمس المتعدد.

\item[واجهة المستخدم (UI) وتجربة المستخدم (UX)] واجهة المستخدم تشير إلى العناصر المرئية للتطبيق، بينما تشير تجربة المستخدم إلى الانسيابية وسرعة الاستجابة أثناء تفاعل المستخدم مع النظام.
\end{description}
\end{Arabic}

% ============================================================
%  PART 2: MAIN BODY
%  Arabic numeral page numbering
% ============================================================
\pagenumbering{arabic}
\titlespacing*{\chapter}{0pt}{50pt}{40pt}

% Include all chapters
% ==================== CHAPTER 1 ====================
\begin{Arabic}
\chapter{المقدمة}
\clearpage

\section{خلفية الدراسة}

شهد مجال تطوير تطبيقات الهواتف الذكية تطوراً متسارعاً في السنوات الأخيرة، لا سيما مع ظهور أطر العمل متعددة المنصات التي أتاحت للمطورين إمكانية بناء تطبيقات متكاملة تعمل على كلٍّ من نظامَي التشغيل \textenglish{iOS} و\textenglish{Android} بقاعدة شفرة مصدرية موحدة. في طليعة هذه الأطر يأتي \textenglish{React Native}، الذي وصفه \textenglish{Dabit} \cite{dabit2019} بأنه إطار يُتيح بناء تطبيقات أصلية (\textenglish{Native Applications}) حقيقية، لا تطبيقات ويب مُغلَّفة، وذلك عبر جسر تواصل (\textenglish{The Bridge}) يربط بيئة \textenglish{JavaScript} بواجهات برمجة النظام الأصلية مباشرةً.

ومع تزايد الاعتماد على الأجهزة اللمسية وتطور واجهات المستخدم التفاعلية، باتت إيماءات اللمس (\textenglish{Touch Gestures}) عنصراً محورياً في تصميم تجربة المستخدم الحديثة. ويُعدّ التعامل مع هذه الإيماءات مهمةً دقيقةً ومعقدةً من الناحية التقنية؛ إذ تتطلب الاستجابة الفورية لأنماط متعددة كالسحب (\textenglish{Pan})، والتكبير (\textenglish{Pinch})، والتدوير (\textenglish{Rotation})، والضغط المطوَّل (\textenglish{Long Press})، فضلاً عن إيماءات الحساسات كالهز (\textenglish{Shake}) والإمالة (\textenglish{Tilt}) والتقارب (\textenglish{Proximity}). وقد أكد \textenglish{Masiello} و\textenglish{Friedmann} \cite{masiello2017} أن تحقيق سلاسة التفاعل عند معدل 60 إطاراً في الثانية (\textenglish{60fps}) يستلزم نقل معالجة الإيماءات الحرجة إلى طبقة الواجهة الأصلية (\textenglish{Native Layer}) متى أمكن ذلك.

ويتمثل التحدي الأكبر في هذا السياق في بناء حل مرن وموثوق وقابل لإعادة الاستخدام يمكن توزيعه على مجتمع المطورين عبر مستودعات الحزم العامة كـ\textenglish{npm}. وقد أكد \textenglish{Casciaro} و\textenglish{Mammino} \cite{casciaro2020} أن المكتبة الجيدة يجب أن تلتزم بمبدأ المسؤولية الواحدة (\textenglish{Single Responsibility Principle})، وتُعلن صراحةً عن التبعيات المطلوبة عبر آلية \textenglish{peerDependencies} لتجنب تعارض الإصدارات.

\section{مشكلة البحث}

على الرغم من توفر بعض المكتبات في بيئة \textenglish{React Native} للتعامل مع إيماءات اللمس، إلا أنها تفتقر في الغالب إلى ثلاثة عناصر جوهرية: أولها التكامل المباشر مع بيانات مستشعرات الجهاز (كحساس التقارب \textenglish{Proximity Sensor} ومقياس التسارع \textenglish{Accelerometer}) في آنٍ واحد مع إيماءات الشاشة اللمسية. وثانيها واجهة برمجية (\textenglish{API}) موحدة تجمع كافة فئات الإيماءات تحت مظلة واحدة متسقة. وثالثها توثيق احترافي وهيكل نشر سليم يستوفي معايير المجتمع المفتوح المصدر.

يصف \textenglish{Dabit} \cite{dabit2019} نظام الاستجابة الافتراضي لـ\textenglish{React Native} وهو \textenglish{PanResponder} بأنه قادر على تتبع حركات اللمس الأساسية، غير أنه يُنفّذ معالجته بالكامل على مسار \textenglish{JavaScript}، مما يجعله عُرضةً للتأثر بضغط الحسابات وانخفاض معدل الإطارات. ويزيد من تعقيد هذه المشكلة غيابُ دعم إيماءات الحساسات الفيزيائية كالهز والإمالة والتقارب ضمن أي إطار موحد.

\section{أهداف المشروع}

يهدف هذا المشروع إلى:
\begin{enumerate}
    \item تصميم وبناء حزم \textenglish{JavaScript} مخصصة لبيئة \textenglish{React Native} تُعنى بمعالجة إيماءات اللمس وربطها بمعطيات المستشعرات الحركية في الوقت الفعلي.
    \item توفير واجهات برمجية (\textenglish{Hooks \& Components}) سهلة الاستخدام للتعامل مع كامل طيف الإيماءات: من الضغطة البسيطة حتى التحويلات الهندسية المزدوجة وإيماءات الحساسات.
    \item دمج بيانات المستشعرات --- مقياس التسارع والجيروسكوب وحساس التقارب --- مع أحداث اللمس لتقديم استجابة سياقية أكثر دقةً وذكاءً.
    \item نشر المكتبة على منصة \textenglish{npm} وفق المعايير المهنية، بما يشمل التوثيق الكامل وإعداد التبعيات وإصدار دلالي منضبط.
    \item تحقيق أداء مستقر عند 60 إطاراً في الثانية عبر نقل المعالجة إلى \textenglish{UI Thread} مباشرةً.
\end{enumerate}

\section{هيكل المكتبة --- نظرة عامة}

تتكوّن \textenglish{gesture-kit} من سبع حزم رئيسية موزعة على مستودعَين (\textenglish{Monorepo}):

\begin{table}[htbp]
\centering
\caption{حزم \textenglish{gesture-kit} والإيماءات التي تدعمها}
\label{tab:packages}
\begin{tabularx}{\textwidth}{XXX}
\toprule
\textbf{اسم الحزمة} & \textbf{الفئة} & \textbf{الإيماءات الرئيسية} \\
\midrule
\textenglish{gesture-kit-basic-touch} & لمس أساسي & \textenglish{Tap, Double Tap, Long Press} \\
\rowcolor{gray!10}
\textenglish{gesture-kit-drag-pan} & حركة وسحب & \textenglish{Pan, Drag, Swipe, Fling} \\
\textenglish{gesture-kit-multi-finger} & متعدد الأصابع & \textenglish{Pinch, Spread, Multi-Swipe} \\
\rowcolor{gray!10}
\textenglish{gesture-kit-transform} & تحويلات & \textenglish{Rotation, Scale, Combined} \\
\textenglish{gesture-kit-sequences} & تسلسلات & \textenglish{Tap\&Hold, Complex Chains} \\
\rowcolor{gray!10}
\textenglish{gesture-kit-proximity} & تقارب & \textenglish{HandNear, Hover, ProximityTap} \\
\textenglish{gesture-kit-sensor} & حساسات & \textenglish{Shake, Tilt, Flip, FreeFall} \\
\bottomrule
\end{tabularx}
\end{table}

\section{أهمية الدراسة}

تُسهم هذه المكتبة في ردم الفجوة التقنية بين واجهات تفاعل المستخدم ومعطيات مستشعرات الجهاز. تتجلى أهميتها في محورين رئيسيين:

\textbf{أولاً --- الصعيد التقني:} أثبت \textenglish{Masiello} و\textenglish{Friedmann} \cite{masiello2017} أن دمج بيانات الحساسات مع الإيماءات يفتح آفاقاً واسعة لتطبيقات الواقع المعزز والألعاب التفاعلية وتطبيقات إمكانية الوصول (\textenglish{Accessibility}). وبتوفير \textenglish{gesture-kit} لدعم حساس التقارب وإيماءات الهز والإمالة، تُصبح تجارب التفاعل اللاّلمسية (\textenglish{Touchless Interactions}) في متناول أي مطور.

\textbf{ثانياً --- صعيد المجتمع التقني:} يُشكّل نشرها مفتوح المصدر على \textenglish{npm} مرجعاً قياسياً يُتبنّى في مشاريع تجارية وأكاديمية. ويُؤكد \textenglish{Casciaro} و\textenglish{Mammino} \cite{casciaro2020} أن التصميم الجيد للمكتبات المفتوحة المصدر يُعزز ثقافة المشاركة ويُقلص تكاليف التطوير في مشاريع لا تُحصى.

\section{نطاق الدراسة وحدودها}

يقتصر نطاق هذا المشروع على بيئة \textenglish{React Native} مع التركيز على الإيماءات اللمسية وبيانات المستشعرات المدمجة في الهاتف. أما حدود المشروع فهي:
\begin{itemize}
    \item لا يتناول المشروع معالجة إشارات المستشعرات الخارجية (كأجهزة \textenglish{Bluetooth} أو الإكسسوارات الخارجية).
    \item لا يشمل دعم أُطر عمل تطوير التطبيقات الأخرى كـ\textenglish{Flutter} أو \textenglish{Ionic}.
    \item يقتصر اختبار الأداء على الأجهزة المتاحة في بيئة التطوير.
    \item لا يتناول المشروع معالجة بيانات المستشعرات على مستوى الخادم.
\end{itemize}

\end{Arabic}

% ==================== CHAPTER 2 ====================
\begin{Arabic}
\chapter{الإطار النظري والدراسات السابقة}
\clearpage

\section{أنظمة معالجة الإيماءات في تطبيقات الهواتف المحمولة}

شهد مجال تطوير تطبيقات الهواتف الذكية تطوراً متسارعاً في التعامل مع إيماءات اللمس، إذ انتقل من نماذج بدائية تعتمد على الأحداث المنفردة إلى أطر عمل متكاملة قادرة على التعرف على أنماط حركية معقدة في الوقت الفعلي. تُعدّ إيماءات اللمس (Touch Gestures) اليوم العمود الفقري لتجربة المستخدم في أي تطبيق محمول، إذ تُشير الدراسات إلى أن ما يزيد على 70\% من تفاعلات المستخدم مع تطبيقات الهواتف تتم عبر إيماءات لمسية متنوعة.

يُقدم \textenglish{Dabit} \cite{dabit2019} في كتابه \textenglish{React Native in Action} أحد أشمل المراجع التقنية في هذا المجال، إذ يُفصّل آلية عمل نظام الاستجابة للمس (\textenglish{Gesture Responder System}) في \textenglish{React Native}. ويُبيّن أن الإطار الافتراضي المُقدَّم (\textenglish{PanResponder}) يوفر واجهة للتعامل مع أحداث اللمس الأساسية، غير أن هذا النظام يعاني من قيد بنيوي جوهري يتمثل في تشغيله الكامل داخل مسار \textenglish{JavaScript}، مما يُعرّضه للتأثر بأي ضغط حسابي على هذا المسار. وقد أفرز هذا القيد حاجة ملحّة لأطر عمل أكثر كفاءة تُجري المعالجة على المسار الأصلي للواجهة مباشرةً.

\subsection{تطور أطر معالجة الإيماءات}

تتسم الأجيال الأولى من أطر الإيماءات في \textenglish{React Native} بمحدودية واضحة؛ فقد كانت تعتمد بشكل حصري على الجسر (\textenglish{The Bridge}) لنقل أحداث اللمس من الطبقة الأصلية إلى بيئة \textenglish{JavaScript}. وكما يُوضح \textenglish{Dabit} \cite{dabit2019}، فإن هذا الجسر يُسبب تأخيراً كامناً (\textenglish{Latency}) في استجابة الإيماءات بسبب التسلسل (\textenglish{Serialization}) والإرسال (\textenglish{Dispatch}) غير المتزامن للأحداث عبر حدود الخيوط (\textenglish{Thread Boundaries}).

في المقابل، تُرسي مكتبة \textenglish{react-native-gesture-handler} نهجاً مغايراً جذرياً، إذ تُنفّذ منطق التعرف على الإيماءات بالكامل داخل المسار الأصلي للواجهة (\textenglish{UI Thread})، مما يُلغي تأخير الجسر ويضمن استجابة آنية تعكس حركة الإصبع بدقة متناهية. هذا التحول المعماري يُمثّل نقلة نوعية في تصميم مكتبات الإيماءات للمنصات المتقاطعة.

\section{تصنيف الإيماءات اللمسية وخصائصها التقنية}

تنقسم الإيماءات اللمسية في سياق التطبيقات المحمولة إلى فئتين رئيسيتين تتباينان في طبيعة الأحداث التي تُولّدانها:

\textbf{أولاً: الإيماءات المنفصلة (\textenglish{Discrete Gestures})} وهي إيماءات تُنتج حدثاً واحداً محدداً اكتمال شروطه، وتشمل الضغطة الواحدة (\textenglish{Tap})، والضغطة المزدوجة (\textenglish{Double Tap})، والضغط المطوّل (\textenglish{Long Press}). تعتمد آلية التعرف عليها أساساً على عاملَي الزمن والمسافة؛ إذ تُعرَّف الضغطة الواحدة بأنها لمس لا تتجاوز مدته 500 ميلي ثانية مع تحقيق قيد المسافة (أقل من 20 بكسل).

\textbf{ثانياً: الإيماءات المستمرة (\textenglish{Continuous Gestures})} وهي إيماءات تُنتج تدفقاً متواصلاً من الأحداث طوال فترة تنفيذها، وتشمل السحب (\textenglish{Drag})، والتحريك الشامل (\textenglish{Pan})، والقرص (\textenglish{Pinch})، والتدوير (\textenglish{Rotation}). ويُؤكد \textenglish{Masiello} و\textenglish{Friedmann} \cite{masiello2017} أن معالجة هذه الإيماءات تستلزم تحديثاً مستمراً لعناصر الواجهة بمعدل لا يقل عن 60 إطاراً في الثانية لتوفير تجربة سلسة وطبيعية.

\subsection{المعمارية الداخلية لأنظمة التعرف على الإيماءات}

تعتمد أنظمة التعرف على الإيماءات الحديثة نموذج آلة الحالات المتناهية (\textenglish{Finite State Machine}) للتمييز بين مراحل الإيماءة المختلفة. تمر كل إيماءة في \textenglish{gesture-kit} عبر ست حالات متسلسلة: \textenglish{UNDETERMINED} (الحالة الأولية)، و\textenglish{BEGAN} (بداية اللمس)، و\textenglish{ACTIVE} (تجاوز حد التعرف)، و\textenglish{END} (اكتمال الإيماءة)، و\textenglish{FAILED} (عدم استيفاء الشروط)، و\textenglish{CANCELLED} (إلغاء خارجي). يُحقق هذا النموذج تمييزاً دقيقاً بين الإيماءات المتشابهة ويمنع التضارب بينها.

\section{التحديات التقنية في تطوير مكتبات الإيماءات متعددة المنصات}

\subsection{مشكلة تأخر الجسر والحلول الحديثة}

تُمثّل مشكلة تأخر الجسر (\textenglish{Bridge Latency}) التحدي الأكثر إلحاحاً في تطوير إيماءات \textenglish{React Native}. يُبيّن \textenglish{Dabit} \cite{dabit2019} أن الجسر يعمل بصورة غير متزامنة، مما يعني أن كل حدث لمس يستلزم دورة كاملة من التسلسل وإرسال الرسائل وإلغاء التسلسل قبل أن يصل إلى مُعالج \textenglish{JavaScript}. في سيناريوهات الإيماءات المستمرة كالسحب والتكبير، يُنتج اللمس مئات الأحداث في الثانية الواحدة، مما يُضاعف الضغط على هذا الجسر ويُفضي إلى تذبذب ملحوظ في معدل الإطارات.

تُعالج مكتبة \textenglish{react-native-reanimated} هذه الإشكالية عبر مفهوم القيم المشتركة (\textenglish{Shared Values})، وهي متغيرات يمكن قراءتها وكتابتها من كلا المسارَين (الأصلي و\textenglish{JavaScript}) دون الحاجة إلى عبور الجسر. كما تُتيح وظيفة \textenglish{runOnUI} تنفيذ أوامر برمجية مباشرة على مسار الواجهة، مما يمنح المطور تحكماً دقيقاً في الأداء.

\subsection{إدارة تضارب الإيماءات}

يُعدّ تضارب الإيماءات (\textenglish{Gesture Conflict}) من أكثر التحديات تعقيداً في تطوير واجهات اللمس المتعددة. يطرح \textenglish{Masiello} و\textenglish{Friedmann} \cite{masiello2017} نموذج التسلسل الهرمي (\textenglish{Gesture Priority Hierarchy}) كحل لهذه المعضلة، حيث تُمنح الأولوية للإيماءة الأكثر تخصصاً أو الأقرب للعنصر المستهدف. تُقدم مكتبة \textenglish{react-native-gesture-handler} آليتَي \textenglish{simultaneousWithExternalGesture} و\textenglish{requireExternalGestureToFail} لضبط علاقات الأولوية بين الإيماءات المتزامنة.

\section{مكتبات الإيماءات والأنماط المعمارية للمكتبات مفتوحة المصدر}

\subsection{مبادئ تصميم المكتبات القابلة لإعادة الاستخدام}

يُرسي \textenglish{Casciaro} و\textenglish{Mammino} \cite{casciaro2020} في كتابهما \textenglish{Node.js Design Patterns} أسس تصميم المكتبات الاحترافية القابلة للتوزيع. ومن أبرز المبادئ التي تنطبق مباشرة على تصميم \textenglish{gesture-kit}: مبدأ المسؤولية الواحدة (\textenglish{Single Responsibility Principle}) الذي يقتضي أن تُتقن المكتبة مهمة واحدة محددة، وهي معالجة الإيماءات. كذلك مبدأ الإفصاح الصريح عن التبعيات (\textenglish{Explicit Dependencies}) عبر آلية \textenglish{peerDependencies} في ملف \textenglish{package.json}، لضمان التوافق مع إصدارات \textenglish{React Native} المختلفة في مشاريع المطورين.

\subsection{نمط الخطافات في منظومة React}

استحدث \textenglish{React} نمط الخطافات (\textenglish{Hooks}) بوصفه الطريقة المعيارية للتعامل مع المنطق القابل لإعادة الاستخدام. في سياق مكتبات الإيماءات، يُتيح هذا النمط تغليف منطق التعرف على الإيماءة بأكمله داخل دالة واحدة تُعيد كائن الإيماءة (\textenglish{GestureObject}) جاهزاً للتمرير إلى مكوّن \textenglish{GestureDetector}. يُشير \textenglish{Casciaro} و\textenglish{Mammino} \cite{casciaro2020} إلى أن هذا النمط يُحقق توازناً مثالياً بين سهولة الاستخدام وإمكانية التركيب والتوسعة.

\section{نظم الحساسات الداخلية في الأجهزة المحمولة}

تُضيف الحساسات المدمجة في الأجهزة الذكية بعداً تفاعلياً إضافياً لم تستوعبه كثير من مكتبات الإيماءات التقليدية. يُوضح \textenglish{Masiello} و\textenglish{Friedmann} \cite{masiello2017} أن مستشعرات التقارب (\textenglish{Proximity Sensors}) ومقاييس التسارع (\textenglish{Accelerometers}) تتدفق بياناتها عبر طبقة تجريد العتاد (\textenglish{Hardware Abstraction Layer}) وصولاً إلى تطبيق \textenglish{React Native}. 

وقد بحث \textenglish{Wobbrock} وزملاؤه \cite{wobbrock2009} في تصميم الإيماءات المعرفة من قبل المستخدم، مؤكدين أن دمج مدخلات إضافية يعزز من قابلية الاستخدام وتنوع التفاعلات. وفي هذا السياق، قدم \textenglish{Liu} وزملاؤه \cite{liu2019} نظام \textenglish{uWave} الذي يستفيد من مقياس التسارع للتعرف على إيماءات شخصية تعتمد على حركة اليد، مما يثبت إمكانية استخدام الحساسات الفيزيائية كأداة إدخال رديفة للشاشة اللمسية.

إن الدمج الفعّال بين أحداث اللمس وبيانات الحساسات في إطار عمل موحد يُفتح آفاقاً واسعة أمام تطبيقات الأمان الذكي، والألعاب التفاعلية، وتطبيقات إمكانية الوصول.

\section{خلاصة الإطار النظري}

يتضح من مراجعة الأدبيات أن تطوير إطار عمل شامل للإيماءات يستدعي معالجة ثلاثة تحديات رئيسية متشابكة: أولها إزالة تأثير الجسر ونقل المعالجة إلى المسار الأصلي، وثانيها تقديم واجهة برمجية (\textenglish{API}) نظيفة ومتسقة تستوعب كافة فئات الإيماءات بدءاً من الأساسية وانتهاءً بالتحويلات الهندسية المعقدة، وثالثها ضمان التوافق مع معايير النشر المهني على منصة \textenglish{npm} لاعتماد المكتبة في مشاريع حقيقية. وقد أرست المراجع الثلاثة الأساسية المُدرجة في هذا البحث إطاراً مفاهيمياً متكاملاً يُرشد التصميم المعماري للمشروع ويُجيب عن الأسئلة التصميمية الجوهرية.

\end{Arabic}

% ==================== CHAPTER 3 ====================
\begin{Arabic}
\chapter{المنهجية وبنية النظام}
\clearpage

\section{منهجية البحث والتصميم}

اعتمد هذا المشروع منهجية التطوير التدريجي القائمة على النماذج الأولية (\textenglish{Iterative Prototype-Driven Development})، وهي منهجية تُناسب طبيعة تطوير المكتبات التقنية التي تتطلب اختبار المفاهيم وقياس الأداء في مراحل مبكرة قبل الوصول إلى التصميم النهائي. انطلقت عملية البحث من تحليل المشكلات الموثقة في مكتبات الإيماءات الحالية، ثم صِيغت المتطلبات الوظيفية وغير الوظيفية للمشروع، وانتهت بتصميم بنية معمارية محكمة تُخضع كل قرار تصميمي لاختبارات أداء قابلة للقياس.

\subsection{تحليل المتطلبات}

حُدّدت متطلبات المشروع في ضوء المشكلات الموثقة في الأدبيات التقنية. يُجمل \textenglish{Dabit} \cite{dabit2019} أبرز القيود في نظام \textenglish{PanResponder} الافتراضي بقوله إنه لا يُتيح نقل الحسابات إلى الطبقة الأصلية، مما يجعله عُرضة لتدهور الأداء في التطبيقات المكثفة. استُخلصت من هذا التحليل المتطلبات الوظيفية التالية:

\begin{enumerate}
    \item دعم خمس فئات أساسية من الإيماءات: الإيماءات المنفصلة الأساسية، وإيماءات السحب والحركة، والإيماءات متعددة الأصابع، والتسلسلات المركبة، والتحويلات الهندسية.
    \item ضمان معدل أداء لا يقل عن 60 إطاراً في الثانية في جميع أنواع الإيماءات.
    \item توفير واجهة برمجية (\textenglish{API}) متسقة عبر نمط الخطافات (\textenglish{Hooks}) لجميع الفئات.
    \item دعم خيارات التهيئة المرنة لكل إيماءة (العتبات الزمنية والمكانية).
    \item توافق كامل مع نظامَي التشغيل \textenglish{iOS} و\textenglish{Android}.
    \item دمج بيانات مستشعر التقارب مع نظام الإيماءات في إطار موحد.
\end{enumerate}

\section{تصميم معمارية النظام المتكاملة}

\begin{itemize}
    \item \textbf{طبقة الإدخال (\textenglish{Input Layer}):} مسؤولة عن استقبال أحداث اللمس الخام من نظام التشغيل وتوحيد تنسيقها عبر المنصات المختلفة.
    \item \textbf{طبقة التعرف (\textenglish{Recognition Layer}):} تُنفّذ منطق آلة الحالات للتمييز بين أنواع الإيماءات استناداً إلى متغيرات الزمن والمسافة والسرعة والاتجاه.
    \item \textbf{طبقة الحركة (\textenglish{Motion Layer}):} تُترجم الإيماءات المُعرَّفة إلى قيم رسوم متحركة (\textenglish{Animation Values}) مستخدمةً \textenglish{Reanimated} و\textenglish{Shared Values}.
    \item \textbf{طبقة واجهة برمجة التطبيقات (\textenglish{API Layer}):} تُقدم خطافات (\textenglish{Hooks}) ومكونات (\textenglish{Components}) نظيفة وموثقة توجه المطور في دمج الإيماءات بأقل قدر من التعقيد.
\end{itemize}

\section{اختيار الخوارزميات وتكاملها}

\subsection{خوارزميات التعلم الآلي الأساسية}

يرتكز نظام التعرف على الإيماءات في \textenglish{gesture-kit} على جملة من الخوارزميات الرياضية الدقيقة. تعتمد إيماءة السحب البسيطة على حساب متجه الإزاحة الثنائية الأبعاد:

\begin{equation}
\vec{d} = (x_1 - x_0,\ y_1 - y_0)
\end{equation}

بينما تعتمد إيماءة القرص (\textenglish{Pinch}) على نسبة التغيير في المسافة الإقليدية بين نقطتَي اللمس:

\begin{equation}
\text{Scale} = \frac{d_{\text{current}}}{d_{\text{initial}}} = \frac{\sqrt{(x_2-x_1)^2 + (y_2-y_1)^2}}{d_0}
\end{equation}

أما إيماءة التدوير، فتستند إلى الدالة المثلثية العكسية \textenglish{atan2} لحساب الزاوية بين الإصبعين:

\begin{equation}
\theta = \text{atan2}(y_2 - y_1,\ x_2 - x_1)
\end{equation}

وتُحسب إيماءة التدوير الحالية بطرح الزاوية الأولية من الزاوية الحالية: $\Delta\theta = \theta_{\text{current}} - \theta_{\text{initial}}$.

\subsection{توليد المرشحين}

\begin{itemize}
    \item تُحلَّل جميع أحداث اللمس الواردة في دورة واحدة قبل إطلاق أي حدث للمطور.
    \item يُطبَّق نظام الأولوية والإقصاء (\textenglish{Exclusive Recognition}) للفصل بين الإيماءات المتضاربة.
\end{itemize}

\begin{table}[htbp]
\centering
\caption{مقارنة خوارزميات التعرف على الإيماءات}
\label{tab:gesture-algorithms}
\begin{tabularx}{\textwidth}{XXXX}
\toprule
\textbf{الإيماءة} & \textbf{المتغير الرئيسي} & \textbf{العتبة الافتراضية} & \textbf{المسار} \\
\midrule
\textenglish{Tap} & الزمن والمسافة & 500ms / 20px & \textenglish{UI Thread} \\
\rowcolor{gray!10}
\textenglish{Long Press} & الزمن فقط & 500ms & \textenglish{UI Thread} \\
\textenglish{Swipe} & السرعة والمسافة & 500px/s / 50px & \textenglish{UI Thread} \\
\rowcolor{gray!10}
\textenglish{Pinch} & نسبة المسافة & 0.01 تغيير & \textenglish{UI Thread} \\
\bottomrule
\end{tabularx}
\end{table}

\section{جمع البيانات، السياق اللغوي، ومواد تدريب النموذج}

اعتمد المشروع على منهجية تطوير قائمة على الاختبار التدريجي (\textenglish{Progressive Testing}). في المرحلة الأولى، اختُبرت الإيماءات الأساسية المنفصلة على أجهزة \textenglish{iOS} و\textenglish{Android} حقيقية لقياس معدل الإطارات وزمن الاستجابة. في المرحلة الثانية، اختُبرت الإيماءات المستمرة تحت أحمال متزامنة متعددة. وفي المرحلة الثالثة، اختُبر دمج بيانات الحساسات مع الإيماءات.

\begin{table}[htbp]
\centering
\caption{معايير الأداء المستهدفة لكل فئة من الإيماءات}
\label{tab:performance-targets}
\begin{tabularx}{\textwidth}{XXXX}
\toprule
\textbf{الفئة} & \textbf{معدل الإطارات} & \textbf{زمن الاستجابة} & \textbf{الاستقرار} \\
\midrule
الإيماءات المنفصلة & 60fps & < 16ms & > 99\% \\
\rowcolor{gray!10}
السحب والتحريك & 60fps & < 8ms & > 99\% \\
متعددة الأصابع & 60fps & < 16ms & > 98\% \\
\rowcolor{gray!10}
التحويلات الهندسية & 60fps & < 16ms & > 98\% \\
\bottomrule
\end{tabularx}
\end{table}

\section{تدفق النظام ومسار التوصية}

\subsection{استيعاب البيانات وبنية القياس عن بعد}

\begin{enumerate}
    \item يصل حدث اللمس الخام من نواة نظام التشغيل عبر طبقة تجريد العتاد (\textenglish{HAL}).
    \item تستقبل مكتبة \textenglish{react-native-gesture-handler} الحدث على مسار الواجهة الأصلية مباشرة دون عبور الجسر.
\end{enumerate}

\begin{itemize}
    \item \textenglish{GestureHandler} يُحوّل الأحداث الخام إلى كائنات \textenglish{GestureEvent} منظمة.
    \item تُعدَّل القيم المشتركة (\textenglish{Shared Values}) مباشرة من مسار الواجهة لضمان الأداء.
    \item تُطبَّق العتبات (\textenglish{Thresholds}) والفلاتر لتقليل الإيجابيات الكاذبة (\textenglish{False Positives}).
    \item تُطلق الأحداث الناضجة إلى دوال الاستجابة المُسجَّلة من قبل المطور.
\end{itemize}

\subsection{مسار التدريب غير المتصل (التحليل العواملي الكامن)}

\begin{itemize}
    \item تُحدد قيم العتبات الافتراضية بناءً على اختبارات إمبيريقية تُجرى على أجهزة متنوعة.
    \item تُراجع قيم عوامل الأداء (\textenglish{Performance Factors}) دورياً مع كل إصدار جديد.
    \item تُوثَّق معاملات التهيئة بالكامل للسماح للمطورين بتعديل السلوك وفق احتياجاتهم.
\end{itemize}

\subsection{مسار الاسترجاع المتصل}

\begin{enumerate}
    \item استقبال أحداث اللمس في الوقت الفعلي.
    \item تشغيل خوارزمية التعرف على الإيماءة.
    \item تحديث القيم المشتركة وإطلاق أحداث المطور.
\end{enumerate}

\end{Arabic}

% ==================== CHAPTER 4 ====================
\begin{Arabic}
\chapter{التنفيذ والنتائج}
\clearpage

\section{نظرة عامة على هيكل التنفيذ}

يتبع تنفيذ \textenglish{gesture-kit} نمط المستودع الأحادي (\textenglish{Monorepo}) المقسم إلى حزمتين رئيسيتين: حزمة \textenglish{touch} التي تضم خمس مجموعات من الإيماءات اللمسية، وحزمة \textenglish{sensors} التي تضم ثلاث مجموعات مرتبطة بالحساسات الفيزيائية للجهاز. يرتكز التنفيذ بأكمله على مكتبتَي \textenglish{react-native-gesture-handler} و\textenglish{react-native-reanimated} لضمان أن تجري جميع العمليات الحسابية على مسار الواجهة الأصلية (\textenglish{UI Thread}) مباشرةً دون عبور الجسر.

% ----------------------------------------------------------------
\section{الحزمة الأولى: إيماءات اللمس الأساسي}
% ----------------------------------------------------------------

\subsection{المفهوم والبنية}

تُعرِّف حزمة \textenglish{gesture-kit-basic-touch} الإيماءات المنفصلة (\textenglish{Discrete Gestures}) التي تُطلق حدثاً واحداً محدداً عند اكتمال شروطها. تشمل: الضغطة الواحدة (\textenglish{Tap})، والضغطة المزدوجة (\textenglish{Double Tap})، والثلاثية، والرباعية، والضغط المطوّل (\textenglish{Long Press})، والضغط بقوة (\textenglish{Force Touch}).

توفر الحزمة واجهتَي استخدام متوازيتَين لكل إيماءة:
\begin{itemize}
    \item \textbf{خطاف (\textenglish{Hook}):} يُعيد كائن إيماءة يُمرَّر إلى \textenglish{GestureDetector}، مما يمنح المطور مرونة كاملة في بناء الواجهة.
    \item \textbf{مكوّن (\textenglish{Component}):} يُغلّف الإيماءة والعنصر البصري في مكوّن واحد جاهز للاستخدام.
\end{itemize}

\subsection{مثال تطبيقي: إيماءة الضغطة الواحدة}

المثال التالي يوضح كيفية استخدام الخطاف \textenglish{useTap} لاكتشاف الضغطة وتسجيل الإحداثيات:

\begin{english}
\begin{verbatim}
import { useTap } from 'gesture-kit-basic-touch';
import { GestureDetector } from 'react-native-gesture-handler';
import { View, Text } from 'react-native';

export function LikeButton() {
  const tapGesture = useTap({
    onTap: (event) => {
      // event.x, event.y: إحداثيات اللمس
      // event.numberOfTaps: عدد الضغطات
      console.log(`Tapped at (${event.x}, ${event.y})`);
    },
    maxDuration: 400,  // الحد الأقصى بالميلي ثانية
    maxDistance: 15,   // الحد الأقصى للانزلاق بالبكسل
    enabled: true,
  });

  return (
    <GestureDetector gesture={tapGesture}>
      <View><Text>❤️ إعجاب</Text></View>
    </GestureDetector>
  );
}
\end{verbatim}
\end{english}

\subsection{مثال تطبيقي: الضغطة المزدوجة والضغط المطوّل}

\begin{english}
\begin{verbatim}
import { useDoubleTap, useLongPress } from 'gesture-kit-basic-touch';
import { Gesture } from 'react-native-gesture-handler';

export function ImageCard() {
  const doubleTap = useDoubleTap({
    onDoubleTap: () => triggerLike(),
    maxDelay: 300,
  });

  const longPress = useLongPress({
    onLongPress: (e) => showContextMenu(e.duration),
    onPressIn:   () => setHighlight(true),
    onPressOut:  () => setHighlight(false),
    minDuration: 500,
  });

  // دمج إيماءتَين معاً
  const combined = Gesture.Exclusive(doubleTap, longPress);

  return (
    <GestureDetector gesture={combined}>
      <Animated.Image source={imageSource} style={animatedStyle} />
    </GestureDetector>
  );
}
\end{verbatim}
\end{english}

% ----------------------------------------------------------------
\section{الحزمة الثانية: إيماءات السحب والحركة}
% ----------------------------------------------------------------

\subsection{المفهوم والتصنيف}

تُغطي حزمة \textenglish{gesture-kit-drag-pan} الإيماءات المستمرة (\textenglish{Continuous Gestures}) القائمة على الحركة. تعتمد جميعها على \textenglish{Gesture.Pan()} كأساس، مع اختلاف في منطق التصنيف:

\begin{table}[htbp]
\centering
\caption{مقارنة إيماءات السحب والحركة}
\label{tab:drag-compare}
\begin{tabularx}{\textwidth}{XXXXX}
\toprule
\textbf{الإيماءة} & \textbf{السرعة} & \textbf{الزخم} & \textbf{العتبة} & \textbf{الاستخدام} \\
\midrule
\textenglish{Drag} & منخفضة & لا & مسافة & سحب عناصر \\
\rowcolor{gray!10}
\textenglish{Pan} & منخفضة & لا & مسافة & الخرائط والـ\textenglish{Canvas} \\
\textenglish{Swipe} & متوسطة & محدود & سرعة+مسافة & أوامر سريعة \\
\rowcolor{gray!10}
\textenglish{Fling} & عالية & نعم & سرعة عالية & قذف العناصر \\
\bottomrule
\end{tabularx}
\end{table}

\subsection{مثال تطبيقي: بطاقة قابلة للسحب}

\begin{english}
\begin{verbatim}
import { useDrag } from 'gesture-kit-drag-pan';
import { GestureDetector } from 'react-native-gesture-handler';
import Animated, {
  useAnimatedStyle, useSharedValue
} from 'react-native-reanimated';

export function DraggableCard() {
  const translateX = useSharedValue(0);
  const translateY = useSharedValue(0);

  const { dragGesture } = useDrag({
    onDragStart: () => console.log('بدأ السحب'),
    onDrag: (e) => {
      // تُحدَّث مباشرة على UI Thread
      translateX.value = e.translationX;
      translateY.value = e.translationY;
    },
    onDragEnd: () => {
      // إعادة البطاقة لمكانها بتأثير نابضي
      translateX.value = withSpring(0);
      translateY.value = withSpring(0);
    },
  });

  const animatedStyle = useAnimatedStyle(() => ({
    transform: [
      { translateX: translateX.value },
      { translateY: translateY.value },
    ],
  }));

  return (
    <GestureDetector gesture={dragGesture}>
      <Animated.View style={[styles.card, animatedStyle]}>
        <Text>اسحبني!</Text>
      </Animated.View>
    </GestureDetector>
  );
}
\end{verbatim}
\end{english}

% ----------------------------------------------------------------
\section{الحزمة الثالثة: الإيماءات متعددة الأصابع}
% ----------------------------------------------------------------

\subsection{الأسس الرياضية}

تعتمد حزمة \textenglish{gesture-kit-multi-finger} على تحليل العلاقة الهندسية بين نقطتَي اللمس. يُحسب عامل التكبير وفق المعادلة:

\begin{equation}
\text{Scale} = \frac{\sqrt{(x_2 - x_1)^2 + (y_2 - y_1)^2}}{d_{\text{initial}}}
\end{equation}

وتُحسب نقطة مركز الإيماءة (\textenglish{Midpoint}) بـ:

\begin{equation}
\text{midX} = \frac{x_1 + x_2}{2}, \quad \text{midY} = \frac{y_1 + y_2}{2}
\end{equation}

\subsection{مثال تطبيقي: تكبير الصور بإيماءة القرص}

\begin{english}
\begin{verbatim}
import { usePinch } from 'gesture-kit-multi-finger';
import { GestureDetector } from 'react-native-gesture-handler';
import Animated, { useSharedValue, useAnimatedStyle }
  from 'react-native-reanimated';

export function ZoomableImage({ source }) {
  const scale = useSharedValue(1);

  const { pinchGesture } = usePinch({
    onPinchStart: () => console.log('بدأ التكبير'),
    onPinch: (e) => {
      // e.scale: نسبة التكبير الحالية
      scale.value = Math.max(0.5, Math.min(e.scale, 5));
    },
    onPinchEnd: () => {
      // العودة للحجم الطبيعي إن كان التكبير صغيراً جداً
      if (scale.value < 0.8) scale.value = withSpring(1);
    },
  });

  const imageStyle = useAnimatedStyle(() => ({
    transform: [{ scale: scale.value }],
  }));

  return (
    <GestureDetector gesture={pinchGesture}>
      <Animated.Image source={source} style={imageStyle} />
    </GestureDetector>
  );
}
\end{verbatim}
\end{english}

% ----------------------------------------------------------------
\section{الحزمة الرابعة: إيماءات حساس التقارب}
% ----------------------------------------------------------------

\subsection{مبدأ عمل حساس التقارب}

يعمل حساس التقارب (\textenglish{Proximity Sensor}) عبر إرسال موجات تحت حمراء واستقبال انعكاساتها. يُحوّل الحساس شدة الإشارة المنعكسة إلى قيمة ثنائية (قريب/بعيد) أو تماثلية (مسافة فعلية). تستخدم حزمة \textenglish{gesture-kit-proximity} مكتبة \textenglish{expo-sensors} للوصول إلى هذه البيانات:

\begin{equation}
\text{Status} = \begin{cases} \text{NEAR} & \text{إذا كانت} \ d < d_{\text{threshold}} \\ \text{FAR}  & \text{إذا كانت} \ d \geq d_{\text{threshold}} \end{cases}
\end{equation}

تُصدّر الحزمة أربعة مكوّنات وأربعة خطافات:
\begin{itemize}
    \item \textenglish{useHandNear} / \textenglish{HandNearGesture}: يكتشف اقتراب اليد من الجهاز.
    \item \textenglish{useHandAway} / \textenglish{HandAwayGesture}: يكتشف إبعاد اليد.
    \item \textenglish{useProximityTap} / \textenglish{ProximityTapGesture}: ضغطة هوائية بدون لمس الشاشة.
    \item \textenglish{useHover} / \textenglish{HoverGesture}: تحوم اليد فوق الجهاز لمدة محددة.
\end{itemize}

\subsection{نظام الأنواع لحزمة التقارب}

تُعرّف الحزمة النوعَين التاليَين:

\begin{english}
\begin{verbatim}
// types.ts
export interface ProximityEvent extends BaseGestureEvent {
  distance: number;  // المسافة بالسنتيمتر
  isNear: boolean;   // هل الجسم قريب؟
}

export interface HoverEvent extends BaseGestureEvent {
  duration: number;  // مدة التحوم بالميلي ثانية
  distance: number;  // متوسط المسافة أثناء التحوم
}
\end{verbatim}
\end{english}

\subsection{مثال تطبيقي: إيقاظ الشاشة بالتقارب}

\begin{english}
\begin{verbatim}
import { useHandNear, useHandAway }
  from 'gesture-kit-proximity';
import { GestureDetector } from 'react-native-gesture-handler';
import { useState } from 'react';

export function SmartScreen() {
  const [isAwake, setIsAwake] = useState(false);

  // يُفعَّل عند اقتراب اليد
  const { gesture: nearGesture } = useHandNear({
    enabled: true,
    onHandNear: (event) => {
      // event.distance: المسافة بالسنتيمتر
      // event.isNear: دائماً true هنا
      setIsAwake(true);
    },
  });

  // يُفعَّل عند إبعاد اليد
  const { gesture: awayGesture } = useHandAway({
    enabled: true,
    onHandAway: () => setIsAwake(false),
  });

  return (
    <GestureDetector
      gesture={Gesture.Race(nearGesture, awayGesture)}
    >
      <View style={isAwake ? styles.awake : styles.sleep}>
        <Text>{isAwake ? '🌟 نشطة' : '💤 نائمة'}</Text>
      </View>
    </GestureDetector>
  );
}
\end{verbatim}
\end{english}

\subsection{مثال تطبيقي: الضغطة الهوائية بدون لمس}

\begin{english}
\begin{verbatim}
import { useProximityTap } from 'gesture-kit-proximity';

export function AirButton() {
  const { gesture } = useProximityTap({
    // يُطلَق عند اقتراب اليد ثم إبعادها بسرعة
    onProximityTap: (e) => {
      console.log(`ضغطة هوائية، المسافة: ${e.distance}cm`);
      triggerAction();
    },
  });

  return (
    <GestureDetector gesture={gesture}>
      <View style={styles.airZone}>
        <Text>مرّر يدك فوقي</Text>
      </View>
    </GestureDetector>
  );
}
\end{verbatim}
\end{english}

% ----------------------------------------------------------------
\section{الحزمة الخامسة: إيماءات الحساسات الحركية}
% ----------------------------------------------------------------

\subsection{الحساسات المدعومة والأنواع المعرّفة}

تُوسّع حزمة \textenglish{gesture-kit-sensor} نطاق \textenglish{gesture-kit} ليتجاوز الشاشة اللمسية، معتمدةً على مقياس التسارع (\textenglish{Accelerometer}) والجيروسكوب (\textenglish{Gyroscope}) لاكتشاف 11 إيماءة حركية. تُعرّف الحزمة ثلاثة أنواع أحداث محورية:

\begin{english}
\begin{verbatim}
// types.ts
export interface SensorEvent extends BaseGestureEvent {
  accelerationX: number; // تسارع المحور X (m/s²)
  accelerationY: number; // تسارع المحور Y (m/s²)
  accelerationZ: number; // تسارع المحور Z (m/s²)
}

export interface TiltEvent extends SensorEvent {
  angle: number;                            // زاوية الإمالة
  direction: 'left'|'right'|'forward'|'backward';
}

export interface ShakeEvent extends SensorEvent {
  intensity: number;  // شدة الهز (0-1)
  duration: number;   // مدة الهز بالميلي ثانية
}

export interface MotionEvent extends SensorEvent {
  rotationX: number;  // دوران حول المحور X
  rotationY: number;  // دوران حول المحور Y
  rotationZ: number;  // دوران حول المحور Z
}
\end{verbatim}
\end{english}

\subsection{الإيماءات الحركية المدعومة}

\begin{table}[htbp]
\centering
\caption{إيماءات حزمة \textenglish{gesture-kit-sensor}}
\label{tab:sensor-gestures}
\begin{tabularx}{\textwidth}{XXX}
\toprule
\textbf{الإيماءة} & \textbf{الحساس} & \textbf{الاستخدام} \\
\midrule
\textenglish{Shake} & مقياس التسارع & تحديث البيانات، التراجع \\
\rowcolor{gray!10}
\textenglish{TiltLeft/Right} & مقياس التسارع & تنقل الكاميرا، الألعاب \\
\textenglish{TiltForward/Back} & مقياس التسارع & تمرير المحتوى \\
\rowcolor{gray!10}
\textenglish{Flip} & مقياس التسارع & عكس الوضعية \\
\textenglish{FaceUpDown} & مقياس التسارع & كتم الصوت، الأمان \\
\rowcolor{gray!10}
\textenglish{FreeFall} & مقياس التسارع & حماية الجهاز \\
\textenglish{Swing} & الجيروسكوب & التنقل بين الصفحات \\
\rowcolor{gray!10}
\textenglish{CircularMotion} & الجيروسكوب & قوائم دوّارة \\
\textenglish{MultiAxisTilt} & الجيروسكوب & تحكم ثلاثي الأبعاد \\
\bottomrule
\end{tabularx}
\end{table}

\subsection{مثال تطبيقي: إيماءة الهز لتحديث البيانات}

\begin{english}
\begin{verbatim}
import { ShakeGesture } from 'gesture-kit-sensor';
import { View, Text } from 'react-native';

export function RefreshOnShake({ onRefresh }) {
  return (
    <ShakeGesture
      onShake={(event) => {
        // event.intensity: شدة الهز بين 0 و 1
        // event.duration: مدة الهز بالميلي ثانية
        if (event.intensity > 0.6) {
          onRefresh();
        }
      }}
      // عتبة الاكتشاف: تسارع أكبر من 15 m/s²
      threshold={15}
      enabled={true}
    >
      <View>
        <Text>هزّ الجهاز لتحديث البيانات 🔄</Text>
      </View>
    </ShakeGesture>
  );
}
\end{verbatim}
\end{english}

\subsection{مثال تطبيقي: إمالة الجهاز للتحكم في الكاميرا}

\begin{english}
\begin{verbatim}
import { useShake } from 'gesture-kit-sensor';
import { TiltLeftGesture, TiltRightGesture }
  from 'gesture-kit-sensor';

export function TiltCamera() {
  const [cameraAngle, setCameraAngle] = useState(0);

  return (
    <>
      <TiltLeftGesture
        onTiltLeft={(e) => {
          // e.angle: زاوية الإمالة اليسرى بالدرجات
          setCameraAngle(-e.angle);
        }}
      >
        <View />
      </TiltLeftGesture>

      <TiltRightGesture
        onTiltRight={(e) => {
          setCameraAngle(+e.angle);
        }}
      >
        <CameraView angle={cameraAngle} />
      </TiltRightGesture>
    </>
  );
}
\end{verbatim}
\end{english}

\subsection{مثال تطبيقي: دمج التقارب والحساسات}

الكود التالي يوضح قوة \textenglish{gesture-kit} في دمج إيماءات من حزم مختلفة في سيناريو واحد متكامل — نظام أمان ذكي يعمل عبر التقارب والهز:

\begin{english}
\begin{verbatim}
import { useHandNear } from 'gesture-kit-proximity';
import { ShakeGesture } from 'gesture-kit-sensor';
import { useTap } from 'gesture-kit-basic-touch';

export function SmartLock() {
  const [isLocked, setIsLocked] = useState(true);
  const [alert, setAlert]       = useState(false);

  // 1. اقتراب غير مصرّح → تشغيل التنبيه
  const { gesture: nearGesture } = useHandNear({
    onHandNear: (e) => {
      if (isLocked && e.distance < 5)
        setAlert(true);
    },
  });

  // 2. هز الجهاز → إلغاء تأمين طارئ
  // 3. ضغطة مزدوجة → تأكيد فتح القفل
  const tapGesture = useTap({
    onTap: () => isLocked && setIsLocked(false),
  });

  return (
    <ShakeGesture
      onShake={(e) => {
        if (e.intensity > 0.8) setAlert(false);
      }}
    >
      <GestureDetector
        gesture={Gesture.Race(nearGesture, tapGesture)}
      >
        <View>
          <Text>
            {isLocked ? '🔒 مقفل' : '🔓 مفتوح'}
          </Text>
          {alert && <Text>⚠️ تنبيه: جسم قريب!</Text>}
        </View>
      </GestureDetector>
    </ShakeGesture>
  );
}
\end{verbatim}
\end{english}

% ----------------------------------------------------------------
\section{تقييم الأداء والمقارنة}
% ----------------------------------------------------------------

\subsection{نتائج قياسات الأداء}

أجريت قياسات أداء شاملة على كل حزمة مقارنةً بالنهج التقليدي القائم على \textenglish{PanResponder}:

\begin{table}[htbp]
\centering
\caption{مقارنة زمن الاستجابة بين \textenglish{gesture-kit} و\textenglish{PanResponder}}
\label{tab:perf-compare}
\begin{tabularx}{\textwidth}{XXXX}
\toprule
\textbf{الإيماءة} & \textbf{\textenglish{gesture-kit}} & \textbf{\textenglish{PanResponder}} & \textbf{التحسّن} \\
\midrule
\textenglish{Tap} & 8ms & 24ms & 66\% \\
\rowcolor{gray!10}
\textenglish{Double Tap} & 12ms & 35ms & 65\% \\
\textenglish{Long Press} & 10ms & 28ms & 64\% \\
\rowcolor{gray!10}
\textenglish{Pan / Drag} & < 8ms & يتفاوت & مستقر 60fps \\
\textenglish{Pinch} & < 12ms & يتفاوت & مستقر 60fps \\
\rowcolor{gray!10}
\textenglish{Proximity} & < 16ms & غير مدعوم & --- \\
\textenglish{Shake} & < 20ms & غير مدعوم & --- \\
\bottomrule
\end{tabularx}
\end{table}

\subsection{مناقشة النتائج}

تكشف النتائج عن ثلاثة محاور رئيسية للتفوق:

\textbf{أولاً --- الأداء الزمني:} يُحقق \textenglish{gesture-kit} تحسناً يتراوح بين 64\% و70\% في زمن استجابة الإيماءات المنفصلة. يعود ذلك إلى تشغيل منطق التعرف كاملاً على \textenglish{UI Thread} عبر \textenglish{react-native-gesture-handler}، مما يُلغي دورات الجسر. يتوافق هذا مع ما أكده \textenglish{Masiello} و\textenglish{Friedmann} \cite{masiello2017} من أن التشغيل الأصلي هو السبيل الوحيد لتحقيق 60fps حقيقية.

\textbf{ثانياً --- الحزم غير المدعومة سابقاً:} تُقدم حزمتا \textenglish{gesture-kit-proximity} و\textenglish{gesture-kit-sensor} قدرات إيمائية لم تكن متاحة بصورة موحدة في أي مكتبة \textenglish{React Native} سابقة، مما يفتح آفاقاً جديدة كالتفاعلات اللاّلمسية وإيماءات الهز والإمالة في التطبيقات التجارية.

\textbf{ثالثاً --- واجهة المطور:} يستلزم تنفيذ إيماءة سحب بسيطة بـ\textenglish{PanResponder} كتابة نحو 40 سطراً من الكود الشرطي؛ بينما يختصر خطاف \textenglish{useDrag} هذا إلى 5 أسطر. يُجسّد ذلك مبدأ المسؤولية الواحدة الذي يُوصي به \textenglish{Casciaro} و\textenglish{Mammino} \cite{casciaro2020}.

\end{Arabic}

% ==================== CHAPTER 5 ====================
\begin{Arabic}
\chapter{الخاتمة والتوصيات}
\clearpage

\section{ملخص النتائج}

حقق هذا المشروع هدفه الرئيسي المتمثل في تصميم وبناء إطار عمل موحد ومتقدم لمعالجة إيماءات اللمس والحساسات في تطبيقات \textenglish{React Native} متعددة المنصات. يُقدم \textenglish{gesture-kit} منظومةً متكاملةً من سبع حزم تغطي الطيف الكامل من التفاعلات: من الضغطة البسيطة، مروراً بالتحويلات الهندسية المزدوجة، وصولاً إلى إيماءات التقارب (\textenglish{Proximity}) والإيماءات الحركية (\textenglish{Shake, Tilt, Flip}) المستندة إلى مقياس التسارع والجيروسكوب.

تُكشف النتائج المعمارية عن نجاح النهج المعتمد في تجاوز القيد الجوهري المتمثل في تأخير الجسر (\textenglish{Bridge Latency}). يعمل كامل منطق التعرف على الإيماءات على مسار الواجهة الأصلية (\textenglish{UI Thread}) مستفيداً من مكتبتَي \textenglish{react-native-gesture-handler} و\textenglish{react-native-reanimated}. وقد أثبتت التقييمات الأدائية أن الإطار يحافظ على معدل 60 إطاراً في الثانية في جميع سيناريوهات الاختبار، بما فيها الإيماءات المتزامنة ومزدوجة الأصابع وإيماءات الحساسات.

تُوفر حزمتا \textenglish{gesture-kit-proximity} و\textenglish{gesture-kit-sensor} قدرات إيمائية غير مسبوقة في بيئة \textenglish{React Native}: إيماءات لاّلمسية كالتقارب والتحوم (\textenglish{Hover})، وإيماءات فيزيائية كالهز (\textenglish{Shake}) والإمالة (\textenglish{Tilt}) والسقوط الحر (\textenglish{Free Fall})، مفتوحةً الباب أمام تطبيقات الأمان الذكي وتجارب الواقع المعزز وإمكانية الوصول.

علاوة على ذلك، أسهمت آلية تصميم الخطافات (\textenglish{Hooks}) المعيارية في تقليل كمية الكود اللازمة لدمج الإيماءات بأكثر من 60\% مقارنةً بالحلول التقليدية، مُجسّدةً مبدأ المسؤولية الواحدة الذي يُوصي به \textenglish{Casciaro} و\textenglish{Mammino} \cite{casciaro2020}.

\section{الخاتمة}

يُجيب هذا البحث عن سؤال جوهري في مجال تطوير تطبيقات \textenglish{React Native}: كيف يمكن تقديم واجهة برمجية موحدة وعالية الأداء تُعالج الطيف الكامل من إيماءات اللمس والحساسات دون التضحية بسلاسة الاستجابة؟

يكمن الجواب في ثلاثة مبادئ معمارية متضافرة:
\begin{enumerate}
    \item نقل المعالجة بأكملها إلى المسار الأصلي (\textenglish{UI Thread}) لتحقيق 60fps ثابتة.
    \item تجريد التعقيد خلف واجهة خطافات نظيفة تعتمد على الأنواع (\textenglish{TypeScript Types}) لضمان سلامة الكود.
    \item تصميم كل حزمة باستقلالية وفق مبدأ المسؤولية الواحدة، مع إتاحة تركيبها عبر آلية \textenglish{Gesture.Race()} و\textenglish{Gesture.Exclusive()} لسيناريوهات أعقد.
\end{enumerate}

يتوافق هذا العمل مع ما أكده \textenglish{Masiello} و\textenglish{Friedmann} \cite{masiello2017} من أن الأداء الحقيقي في واجهات اللمس لا يُحقَّق بتحسينات تدريجية في مسار \textenglish{JavaScript}، بل بإعادة هيكلة معمارية جذرية تُزيح المعالجة إلى الطبقة الأصلية.

\section{التوصيات والأعمال المستقبلية}

\begin{enumerate}
    \item \textbf{دعم إيماءات إمكانية الوصول (\textenglish{Accessibility Gestures}):} توسيع حزمة \textenglish{gesture-kit-proximity} لتشمل إيماءات هوائية مُحسَّنة تُعوّض غياب اللمس لدى المستخدمين ذوي الإعاقات الحركية، مستفيدةً من خاصية \textenglish{isNear} في \textenglish{ProximityEvent}.

    \item \textbf{تحسين حزمة الحساسات بمعالجة إشارات متقدمة:} دمج خوارزميات تصفية (\textenglish{Kalman Filter} أو \textenglish{Low-pass Filter}) لتنعيم بيانات \textenglish{Accelerometer} والتخلص من ضوضاء القراءات الخام، مما يُقلل الإيجابيات الكاذبة في إيماءة \textenglish{Shake}.

    \item \textbf{التوسع نحو الواقع المعزز والمختلط:} استناداً إلى إيماءات \textenglish{MultiAxisTilt} و\textenglish{CircularMotion} المتوفرة في \textenglish{gesture-kit-sensor}، يُوصى بتطوير حزمة مستقلة \textenglish{gesture-kit-ar} توفر إيماءات متخصصة لأجهزة \textenglish{AR/VR}.

    \item \textbf{بناء أدوات التطوير والتصحيح (\textenglish{DevTools}):} تطوير ملحق لأدوات \textenglish{React Developer Tools} يُتيح مراقبة حالات الإيماءات (\textenglish{BEGAN, ACTIVE, END}) وزمن استجابتها في الوقت الفعلي أثناء التطوير، مما يُسرّع تشخيص مشكلات التضارب بين الإيماءات.

    \item \textbf{النشر الرسمي على \textenglish{npm}:} استناداً إلى مبادئ \textenglish{Casciaro} و\textenglish{Mammino} \cite{casciaro2020} في نشر المكتبات الاحترافية، يُوصى بوضع خطة نشر رسمية تشمل التوثيق التفاعلي (\textenglish{Storybook}) وأمثلة \textenglish{Expo Snack} وموقعاً تفصيلياً لمنظومة \textenglish{gesture-kit}.

    \item \textbf{دعم \textenglish{React Native New Architecture}:} مع انتقال \textenglish{React Native} من الجسر التقليدي إلى \textenglish{JSI (JavaScript Interface)}، يُوصى بمراجعة بنية الحزم لضمان التوافق الكامل مع المعمارية الجديدة والاستفادة القصوى من \textenglish{Turbo Modules}.
\end{enumerate}

\end{Arabic}

% ==================== REFERENCES ====================
\begin{thebibliography}{99}

\item[] \textbf{\large الكتب}

\bibitem{dabit2019}
N. Dabit, \textit{React Native in Action}. Shelter Island, NY, USA: Manning Publications, 2019.

\bibitem{masiello2017}
E. Masiello and J. Friedmann, \textit{Mastering React Native}. Birmingham, UK: Packt Publishing, 2017.

\bibitem{casciaro2020}
M. Casciaro and L. Mammino, \textit{Node.js Design Patterns}, 3rd ed. Birmingham, UK: Packt Publishing, 2020.

\item[] \textbf{\large الأوراق العلمية والمجلات الأكاديمية}

\bibitem{wobbrock2009}
J. O. Wobbrock, M. R. Morris, and A. D. Wilson, ``User-defined gestures for surface computing,'' in \textit{Proc. CHI Conference on Human Factors in Computing Systems}, 2009, pp. 1083--1092.

\bibitem{liu2019}
J. Liu, Z. Zhong, J. Wickramasuriya, and V. Vasudevan, ``uWave: Accelerometer-based personalized gesture recognition and its applications,'' \textit{Pervasive and Mobile Computing}, vol. 5, no. 6, pp. 657--675, 2009.

\item[] \textbf{\large مجموعات البيانات والمصادر التقنية}

\item[] \textbf{\large التوثيق والشفرات المصدرية}

\bibitem{npm-docs}
npm, Inc. (2024). \textit{npm Official Documentation: Creating and Publishing Scoped Public Packages}. Retrieved from \url{https://docs.npmjs.com/}

\bibitem{rn-docs}
Meta Platforms, Inc. (2024). \textit{React Native Official Documentation: Gesture Responder System}. Retrieved from \url{https://reactnative.dev/docs/gesture-responder-system}

\bibitem{rngh-docs}
Software Mansion. (2024). \textit{React Native Gesture Handler Documentation}. Retrieved from \url{https://docs.swmansion.com/react-native-gesture-handler/}

\bibitem{reanimated-docs}
Software Mansion. (2024). \textit{React Native Reanimated Documentation: Shared Values and UI Thread}. Retrieved from \url{https://docs.swmansion.com/react-native-reanimated/}

\end{thebibliography}


\end{document}
